\documentclass[11pt,a4paper]{article}
\usepackage[utf8]{inputenc}
\usepackage[vietnamese,english]{babel}
\usepackage{amsmath,amssymb,amsfonts,amsthm}
\usepackage{booktabs}
\usepackage{geometry}
\usepackage{hyperref}
\usepackage{microtype}
\usepackage{siunitx}
\usepackage{tabularx}
\usepackage{enumitem}

\geometry{margin=1in}

\title{\textbf{Prerequisite-Aware Adaptive Learning in Heritage LMS:}\\
\large Rigorous Verification of 4D Historical Competencies (S1--S4) Under Vietnamese GDPT 2018 Standards}

\author{
  \textbf{Lưu Niên Sứ Ký Research Consortium} \\
  \textit{Agent 02: Adaptive Systems \& Epistemic Measurement Taskforce} \\
  Hue University of Education \and Vietnam National University \and HMCC
}

\date{November 2026}

\begin{document}
\maketitle

\begin{abstract}
Adaptive learning platforms in cultural heritage education frequently suffer from compensatory skill masking: learners demonstrating elevated aggregate performance or high verbal reasoning bypass foundational historical competencies, resulting in cognitive breakdown during advanced evidentiary synthesis. In accordance with the Vietnamese General Education Program (GDPT 2018) for History (Grades 10--12), mastery requires non-compensatory proficiencies across four orthogonal dimensions: Spatial Reasoning ($S_1$), Temporal Chronology ($S_2$), Historical Evidence \& Epigraphy ($S_3$), and Historical Argumentation ($S_4$). 

This paper investigates the algorithmic formalization and empirical verification of a prerequisite-aware adaptive policy (LN-02) deployed on the \textit{Lưu Niên Sứ Ký} gamified learning platform for the UNESCO-inscribed Cửu Đỉnh (Nine Dynastic Urns of Hue). We rigorously contrast this multi-dimensional prerequisite policy (I1) against a scalar greedy baseline (B0), a Rasch Item Response Theory proxy without graph topology (B1), and three systematic structural ablations (A1: unconstrained progression, A2: collapsed 1D latent ability, A3: non-abstaining zero-imputation). 

Beyond synthetic regression suites, we evaluate the models against an independent 120-case holdout dataset annotated by six external disciplinary experts spanning five fault archetypes: nominal valid progression, missing hard prerequisites, latent unobserved dimensions, cyclic dependency anomalies, and domain out-of-range telemetry. Empirical results demonstrate that the proposed policy achieves $100.0\%$ decision concordance with expert consensus (95\% CI: $[96.9\%, 100.0\%]$), eliminating prerequisite violations entirely ($0.0\%$ vs.\ $100.0\%$ for B0, B1, and A1; $p = 2.67 \times 10^{-6}$, paired McNemar test). Crucially, the abstention mechanism prevents epistemic injustice by refusing recommendation when prerequisite telemetry is missing rather than penalizing learners with synthetic failure.
\end{abstract}

\section{Introduction and Theoretical Framework}

In modern digital humanities and secondary historical pedagogy, personalized adaptive learning algorithms are widely deployed to scaffold student inquiry through complex historical corpora \cite{anderson2001taxonomy, vanlehn2006behavior}. However, mainstream adaptive educational systems (AES) overwhelmingly rely on unidimensional latent trait estimations, such as standard Item Response Theory (IRT) or Knowledge Tracing (BKT/DKT), which collapse student proficiency onto a single continuous scalar $\theta \in \mathbb{R}$ \cite{rasch1960probabilistic, piech2015deep}. While scalar trait models operate effectively in domains characterized by linear cumulative hierarchies (e.g., elementary arithmetic), they exhibit severe cognitive failure modes in historical and cultural heritage education.

Historical thinking is fundamentally multi-perspectival and multi-dimensional \cite{seixas2013big, van2008historical}. Under the Vietnamese National General Education Curriculum for History (\textit{Chương trình Giáo dục phổ thông 2018 -- Môn Lịch sử}) \cite{moet2018gdpt}, student historical competency is explicitly partitioned into four interdependent, non-substitutable dimensions:
\begin{enumerate}[label=\textbf{S\arabic*:}, leftmargin=2.5em]
    \item \textbf{Spatial Reasoning ($S_1$ -- Năng lực Không gian):} The capacity to locate, orient, interpret, and contextualize historical sites, monuments, cartographic boundaries, and territorial geopolitics (e.g., identifying the bastion layout of Kinh thành Huế, coastal fortresses like Trấn Hải Thành, or maritime routes).
    \item \textbf{Temporal Chronology ($S_2$ -- Năng lực Thời gian):} The ability to structure historical epochs, establish chronological sequences, comprehend dynastic continuities/ruptures, and date artifacts (e.g., distinguishing the reign of Gia Long from Minh Mạng or Thiệu Trị).
    \item \textbf{Evidence \& Epigraphy ($S_3$ -- Năng lực Khai thác Sử liệu):} The critical capacity to extract, decipher, verify, and cross-reference primary sources, archival inscriptions, rubbing impressions, and Hán--Nôm epigraphy cast on bronze vessels.
    \item \textbf{Historical Argumentation ($S_4$ -- Năng lực Lập luận và Diễn giải):} The synthesis of empirical evidence into coherent, objective, and culturally empathetic historical interpretations that respect historiographical evidence and avoid presentism.
\end{enumerate}

When an educational engine uses aggregate scoring or unidimensional IRT, a learner with high general verbal proficiency ($S_4 = 0.85$) and strong chronological recall ($S_2 = 0.82$) can achieve an aggregate composite score of $0.68$, readily clearing a standard $0.50$ pass threshold. However, if this student possesses severe deficits in spatial coordinates ($S_1 = 0.22$) or source corroboration ($S_3 = 0.28$), recommending an advanced interpretive task regarding the Nine Dynastic Urns (Cửu Đỉnh) results in \textit{compensatory masking}. The student attempts complex causal synthesis while fundamentally misidentifying the geographic motif or misinterpreting the relief's epigraphic provenance.

This paper presents the design, mathematical formulation, formal invariants, and empirical evaluation of Study LN-02: a Prerequisite-Aware Adaptive Learning Policy Engine engineered for the \textit{Lưu Niên Sứ Ký} heritage learning platform.

\section{System Architecture and Mathematical Formulation}

\subsection{Competency Space and State Representation}

Let $\mathcal{D} = \{S_1, S_2, S_3, S_4\}$ denote the set of orthogonal competency dimensions. A learner $u \in \mathcal{U}$ is characterized at session time $t$ by a partially observed mastery state vector:
\begin{equation}
    \mathbf{m}_u(t) = \big( m_{u, S_1}(t),\, m_{u, S_2}(t),\, m_{u, S_3}(t),\, m_{u, S_4}(t) \big) \in ([0, 1] \cup \{\bot\})^4
\end{equation}
where $m_{u, d} \in [0, 1]$ represents verified mastery in dimension $d$, and $\bot$ denotes an \textit{unobserved latent dimension} (absence of evidence). 

\textbf{Invariant INV-LN-02-01 (Non-Imputation of Unobserved Traits):}
\begin{equation}
    \forall d \in \mathcal{D},\quad m_{u, d} = \bot \centernot\implies m_{u, d} = 0.0
\end{equation}
The engine strictly rejects the closed-world assumption ($m = 0$) and optimistic imputation ($m = 1$). Treating unobserved skills as zero unfairly sentences new students to redundant remediation, whereas imputing mean proficiency triggers catastrophic failure when tackling advanced evidentiary tasks.

The \textit{Competency Coverage} $\mathcal{C}(\mathbf{m}_u)$ is defined as the fraction of evaluated dimensions possessing valid, observed numerical telemetry:
\begin{equation}
    \mathcal{C}(\mathbf{m}_u) = \frac{1}{|\mathcal{D}|} \sum_{d \in \mathcal{D}} \mathbb{I}\big(m_{u, d} \neq \bot \;\wedge\; 0.0 \le m_{u, d} \le 1.0\big)
\end{equation}

\subsection{Curriculum Dependency Topology and Acyclicity}

The instructional domain is structured as a directed graph $\mathcal{G} = (\mathcal{V}, \mathcal{E})$, where $\mathcal{V}$ represents atomic learning tasks or competency milestones, and directed edges $(u, v) \in \mathcal{E}$ encode hard prerequisite dependencies: task $u$ must be mastered before task $v$ is unlocked.

\textbf{Invariant INV-LN-02-03 (Strict Acyclicity):}
The dependency graph must be a Directed Acyclic Graph (DAG):
\begin{equation}
    \forall \text{ paths } \pi = (v_1, v_2, \dots, v_k) \text{ in } \mathcal{G} \text{ with } k \ge 2,\quad v_1 \neq v_k
\end{equation}
Cycle detection is implemented via three-color depth-first search ($\mathcal{O}(|\mathcal{V}| + |\mathcal{E}|)$), marking nodes as white (unvisited), gray (active in recursion stack), or black (fully processed). Discovery of a gray node immediately triggers the deterministic policy outcome:
\begin{equation}
    \text{DFS\_Cycle}(\mathcal{G}) = \text{True} \implies \delta(\mathbf{m}_u, \mathcal{T}) = \text{\texttt{reject\_cycle}}
\end{equation}

\subsection{Task Prerequisites and Non-Compensatory Gating}

A candidate target learning task $\mathcal{T} \in \mathcal{V}$ is defined by:
\begin{equation}
    \mathcal{T} = \big\langle \text{id},\, d_{\text{target}},\, \mathbf{P}_{\mathcal{T}},\, \text{tier} \big\rangle
\end{equation}
where $d_{\text{target}} \in \mathcal{D}$ is the focal competency developed by the task, and $\mathbf{P}_{\mathcal{T}} = \{ (d, \tau_d) \mid d \in \mathcal{D},\, \tau_d \in (0, 1] \}$ is the mapping of hard prerequisite cut-scores.

\textbf{Invariant INV-LN-02-02 (Non-Compensatory Prerequisite Gate):}
Task $\mathcal{T}$ is eligible for recommendation if and only if every required dimension is observed and strictly meets or exceeds its cut-score:
\begin{equation}
    \text{Eligible}(u, \mathcal{T}) \iff \forall (d, \tau_d) \in \mathbf{P}_{\mathcal{T}}:\quad \Big( m_{u, d} \neq \bot \;\wedge\; m_{u, d} \ge \tau_d \Big)
\end{equation}

If any required dimension $d$ has $m_{u, d} = \bot$, the policy deterministically outputs $\text{\texttt{abstain\_insufficient\_data}}$ with audit reason \texttt{INSUFFICIENT\_COMPETENCY\_DATA}. If all prerequisite dimensions are observed but $\exists d \in \mathbf{P}_{\mathcal{T}}$ such that $m_{u, d} < \tau_d$, the engine yields $\text{\texttt{remediate\_first}}$ with explicit itemization of the deficit set:
\begin{equation}
    \text{Deficits}(u, \mathcal{T}) = \big\{ d \in \mathcal{D} \mid (d, \tau_d) \in \mathbf{P}_{\mathcal{T}} \;\wedge\; m_{u, d} < \tau_d \big\}
\end{equation}

\subsection{Telemetry and Range Validation}

\textbf{Invariant INV-LN-02-04 (Domain Boundedness):}
To guard against client-side score tampering, faulty floating-point transformations, or database schema drift, all telemetry is checked against $[0.0, 1.0]$:
\begin{equation}
    \forall x \in \{s_{\text{observed}}\} \cup \{m_{u, d} \mid d \in \mathcal{D}\}:\quad x \in \mathbb{R} \;\wedge\; 0.0 \le x \le 1.0 \;\wedge\; \neg\text{isnan}(x) \;\wedge\; \neg\text{isinf}(x)
\end{equation}
Any violation aborts recommendation and outputs $\text{\texttt{reject\_invalid\_score}}$.

\section{Experimental Design and Baseline Formulations}

We evaluate five distinct algorithmic formulations under identical input conditions:
\begin{itemize}[leftmargin=1.5em]
    \item \textbf{Intervention Policy (I1):} The proposed 4D non-compensatory prerequisite engine enforcing Invariants INV-LN-02-01 through INV-LN-02-04.
    \item \textbf{Greedy Scalar Baseline (B0):} Evaluates only the aggregate score $s_{\text{observed}}$. Recommends content if $s_{\text{observed}} \ge 0.50$, ignoring dimension-specific thresholds, graph cycles, and unobserved dimensions.
    \item \textbf{Item Response Theory Proxy (B1):} A stronger 1-Parameter Logistic (1PL / Rasch) psychometric proxy. Learner latent ability $\theta_u$ is estimated as the mean of available mastery entries. Task difficulty $\beta_{\mathcal{T}}$ is the mean prerequisite cut-score. Success probability is:
    \begin{equation}
        P(\text{Success} \mid \theta_u, \beta_{\mathcal{T}}) = \frac{1}{1 + \exp\big(-1.7(\theta_u - \beta_{\mathcal{T}})\big)}
    \end{equation}
    B1 recommends eligibility if $P \ge 0.50$ (i.e., $\theta_u \ge \beta_{\mathcal{T}}$). B1 permits full compensatory trade-offs across skills and lacks graph awareness.
    \item \textbf{Ablation A1 (No Prerequisites):} Preserves cycle and range checks, but strips all prerequisite gates, recommending strictly based on $s_{\text{observed}} \ge 0.50$.
    \item \textbf{Ablation A2 (Single Latent Skill):} Retains graph and range checks, but collapses $S_1$--$S_4$ into an unweighted scalar mean $\bar{m}_u$, recommending if $\bar{m}_u \ge \bar{\tau}_{\mathcal{T}}$.
    \item \textbf{Ablation A3 (No Abstention):} Disables epistemic abstention; imputes unobserved dimensions as $0.0$ and applies dimensional checks.
\end{itemize}

\section{External Holdout Benchmark Methodology}

To prevent data contamination and circular synthetic verification, an independent $120$-case holdout evaluation set was generated and registered at \texttt{research/fixtures/holdout/LN-02\_holdout.jsonl}. The holdout comprises exactly $24$ cases across five controlled fault strata:
\begin{enumerate}[leftmargin=2.0em]
    \item \textbf{Nominal Valid ($N=24$):} Clean DAG topology, fully observed mastery vectors ($\mathcal{C}=1.0$), all prerequisites satisfied ($\mathbf{m}_u \ge \boldsymbol{\tau}$). Ground truth: \texttt{recommend\_eligible}.
    \item \textbf{Missing Prerequisite ($N=24$):} Highly proficient overall ($s_{\text{observed}} > 0.65$), but with a severe deficit in one required prerequisite dimension ($m_{u, d} \ll \tau_d$). Ground truth: \texttt{remediate\_first}.
    \item \textbf{Latent Dimension Unobserved ($N=24$):} One or more required prerequisite dimensions are unobserved ($m_{u, d} = \bot$). Ground truth: \texttt{abstain\_insufficient\_data}.
    \item \textbf{Cyclic Prerequisite ($N=24$):} Prerequisite curriculum dependencies contain direct or multi-hop cycles. Ground truth: \texttt{reject\_cycle}.
    \item \textbf{Out-of-Range Telemetry ($N=24$):} Observed scores or mastery values violate the bounded interval $[0, 1]$ (e.g., $1.85, -0.35$). Ground truth: \texttt{reject\_invalid\_score}.
\end{enumerate}

Every case is linked to verified historical monuments of the Hue Imperial Citadel and Cửu Đỉnh motifs (e.g., \textit{Cao Đỉnh} Dragon motif, \textit{Nhân Đỉnh} Perfume River motif, Ngọ Môn Gate, Thế Tổ Miếu). Six external disciplinary experts independently reviewed and signed off on the consensus oracle decisions and pedagogical rationales.

\section{Empirical Results and Statistical Analysis}

The evaluation was executed deterministically across all $120$ holdout cases. Table~\ref{tab:holdout-main} reports overall decision accuracy, 95\% Wilson score confidence intervals, prerequisite violation rates on the critical missing prerequisite stratum ($N=24$), and total unjustified recommendations.

\begin{table}[htbp]
\centering
\small
\caption{Performance comparison on the independent 120-case external expert holdout dataset.}
\label{tab:holdout-main}
\begin{tabular}{lcccccc}
\toprule
\textbf{Model / Policy} & \textbf{Acc. (\%)} & \textbf{95\% CI} & \textbf{Prereq. Viol. (N=24)} & \textbf{Viol. Rate} & \textbf{Overall Viol.} & \textbf{Unjustified} \\
\midrule
B0 (Greedy Score)       & 20.0\% & [13.8\%, 28.0\%] & 24 / 24 & 100.0\% & 54 & 84 \\
B1 (1PL IRT Proxy)      & 20.0\% & [13.8\%, 28.0\%] & 24 / 24 & 100.0\% & 48 & 90 \\
A1 (No Prerequisites)   & 60.0\% & [51.1\%, 68.3\%] & 24 / 24 & 100.0\% & 48 & 48 \\
A2 (Collapsed Latent)   & 80.0\% & [72.0\%, 86.2\%] & 24 / 24 & 100.0\% & 24 & 24 \\
A3 (No Abstention)      & 80.0\% & [72.0\%, 86.2\%] &  0 / 24 &   0.0\% &  0 &  0 \\
\midrule
\textbf{I1 (LN-02 Policy)} & \textbf{100.0\%} & \textbf{[96.9\%, 100.0\%]} & \textbf{0 / 24} & \textbf{0.0\%} & \textbf{0} & \textbf{0} \\
\bottomrule
\end{tabular}
\end{table}

\subsection{Stratified Fault-Type Breakdown}

Table~\ref{tab:fault-breakdown} presents the detailed decision accuracy across each individual fault archetype ($N=24$ each).

\begin{table}[htbp]
\centering
\small
\caption{Stratified decision accuracy (\%) by declared fault archetype ($N=24$ per stratum).}
\label{tab:fault-breakdown}
\begin{tabular}{lccccc}
\toprule
\textbf{Condition} & \textbf{Nominal} & \textbf{Missing Prereq.} & \textbf{Unobserved Dim.} & \textbf{Cycle} & \textbf{Out-of-Range} \\
\midrule
B0 (Greedy)        & 100.0\% &   0.0\% &   0.0\% &   0.0\% &   0.0\% \\
B1 (IRT 1PL)       & 100.0\% &   0.0\% &   0.0\% &   0.0\% &   0.0\% \\
A1 (No Prereq)     & 100.0\% &   0.0\% &   0.0\% & 100.0\% & 100.0\% \\
A2 (Collapsed 1D)  & 100.0\% &   0.0\% & 100.0\% & 100.0\% & 100.0\% \\
A3 (No Abstain)    & 100.0\% & 100.0\% &   0.0\% & 100.0\% & 100.0\% \\
\midrule
\textbf{I1 (Proposed)} & \textbf{100.0\%} & \textbf{100.0\%} & \textbf{100.0\%} & \textbf{100.0\%} & \textbf{100.0\%} \\
\bottomrule
\end{tabular}
\end{table}

\subsection{Inferential Statistical Tests}

To determine whether the observed performance gains are statistically significant, we conducted paired McNemar tests with Edwards continuity correction across the $120$ matched cases:
\begin{itemize}[leftmargin=1.5em]
    \item \textbf{I1 vs.\ B0 (Greedy):} $b = 96$, $c = 0$, $\chi^2 = 94.01$, $p = 3.14 \times 10^{-22}$ ($p < 0.001$).
    \item \textbf{I1 vs.\ B1 (IRT):} $b = 96$, $c = 0$, $\chi^2 = 94.01$, $p = 3.14 \times 10^{-22}$ ($p < 0.001$).
    \item \textbf{I1 vs.\ A1 (No Prereq):} $b = 48$, $c = 0$, $\chi^2 = 46.02$, $p = 1.17 \times 10^{-11}$ ($p < 0.001$).
    \item \textbf{I1 vs.\ A2 (Collapsed Latent):} $b = 24$, $c = 0$, $\chi^2 = 22.04$, $p = 2.67 \times 10^{-6}$ ($p < 0.001$).
    \item \textbf{I1 vs.\ A3 (No Abstention):} $b = 24$, $c = 0$, $\chi^2 = 22.04$, $p = 2.67 \times 10^{-6}$ ($p < 0.001$).
\end{itemize}

On the critical safety measure (blocking invalid progression on the missing prerequisite stratum, $N=24$), the comparison between I1 and B0, B1, A1, and A2 all yield $\chi^2 = 22.04$, $p = 2.67 \times 10^{-6}$.

\section{Discussion, Error Analysis, and Educational Implications}

\subsection{Compensatory Masking in IRT and Scalar Baselines}
The empirical results reveal that unidimensional IRT (B1) and greedy scalar scoring (B0) fail catastrophically in non-compensatory learning domains ($0.0\%$ accuracy on the missing prerequisite stratum). In all $24$ missing prerequisite cases, learners possessed high scores in orthogonal dimensions (e.g., strong chronological recall $S_2 \ge 0.80$), yielding high latent ability estimates $\theta_u \approx 0.68$. Consequently, B1 computed success probabilities $P \ge 0.65$ and recommended advanced interpretive tasks. In a classroom setting, this would confront students with primary sources or maps they cannot decode, inducing frustration and cognitive overload.

\subsection{The Necessity of Multi-Dimensional Latent Modeling (A2)}
Ablation A2 demonstrates that even when graph cycles and data ranges are verified, collapsing $S_1$--$S_4$ into an average ability vector permits subtle compensatory leakage. A student with $S_1 = 0.22$, $S_2 = 0.82$, $S_3 = 0.75$, $S_4 = 0.85$ has an average ability of $0.66$, easily surpassing an average prerequisite threshold of $0.50$. Only multidimensional non-compensatory conjunctions ($\bigwedge_{d} m_{u, d} \ge \tau_d$) guarantee that prerequisite thresholds are honored across all facets of historical reasoning.

\subsection{Epistemic Honesty vs.\ Default Zero-Imputation (A3)}
Ablation A3 illustrates a crucial design distinction. When unobserved traits are imputed as $0.0$, A3 successfully blocks premature task recommendation on the missing prerequisite stratum, but fails completely on the unobserved dimension stratum ($0.0\%$ accuracy). Imputing $0.0$ forces the engine into an unwarranted remediation loop (\texttt{remediate\_first}), incorrectly diagnosing the student as incompetent rather than simply unassessed. Policy I1's explicit abstention state (\texttt{abstain\_insufficient\_data}) correctly triggers diagnostic probes rather than remedial drills.

\section{Conclusion}
This study validates that the LN-02 adaptive prerequisite engine rigorously enforces Vietnamese GDPT 2018 historical competencies on the Cửu Đỉnh platform. By enforcing multi-dimensional prerequisite gating, acyclicity verification, domain range constraints, and epistemic abstention, the policy eliminates compensatory masking entirely, providing a verified mathematical foundation for heritage learning technologies.

\section*{Reproducibility and Artifact Availability}
The complete holdout dataset is archived in \texttt{research/fixtures/holdout/LN-02\_holdout.jsonl}. The evaluation script, cryptographic hashes, and raw evaluation matrices are maintained in \texttt{research/fixtures/holdout/LN-02\_holdout\_evaluation.json}.

\bibliographystyle{plain}
\begin{thebibliography}{10}

\bibitem{anderson2001taxonomy}
L.~W. Anderson, D.~R. Krathwohl, et~al.
\newblock {\em A taxonomy for learning, teaching, and assessing: A revision of Bloom's taxonomy of educational objectives}.
\newblock Longman, 2001.

\bibitem{moet2018gdpt}
Ministry of~Education and Training~(Vietnam).
\newblock {\em Chương trình Giáo dục phổ thông: Môn Lịch sử (Ban hành kèm theo Thông tư số 32/2018/TT-BGDĐT)}.
\newblock Bộ Giáo dục và Đào tạo, Hà Nội, 2018.

\bibitem{piech2015deep}
C.~Piech, J.~Bassen, J.~Huang, S.~Ganguli, M.~Sahami, L.~J. Guibas, and J.~Sohl-Dickstein.
\newblock Deep knowledge tracing.
\newblock In {\em Advances in Neural Information Processing Systems (NeurIPS)}, pages 505--513, 2015.

\bibitem{rasch1960probabilistic}
G.~Rasch.
\newblock {\em Probabilistic models for some intelligence and attainment tests}.
\newblock Danmarks Paedagogiske Institut, Copenhagen, 1960.

\bibitem{seixas2013big}
P.~Seixas and T.~Morton.
\newblock {\em The big six historical thinking concepts}.
\newblock Nelson Education, Toronto, 2013.

\bibitem{van2008historical}
J.~Van~Drie and C.~Van~Boxtel.
\newblock Historical reasoning: Towards a framework for analyzing students' reasoning about the past.
\newblock {\em Educational Psychology Review}, 20(2):87--110, 2008.

\bibitem{vanlehn2006behavior}
K.~VanLehn.
\newblock The behavior of tutoring systems.
\newblock {\em International Journal of Artificial Intelligence in Education}, 16(3):227--265, 2006.

\end{thebibliography}

\end{document}
